% 本文件由 LaTeX 排版 Agent 根据有效 translation.json 自动生成。
% author=Apocrypha; work_id=0813; title=阿黎玛特雅人若瑟的叙述

\chapter{阿黎玛特雅人若瑟的叙述}

% structure: p0001 source=h1 target=omitted reason=page-title-already-rendered-by-parent

\Emph{阿黎玛特雅人若瑟的叙述，他求得了主的身体；其中也述及两个强盗的案情。}\par

% structure: p0003 source=h2 target=section reason=relative-heading-rank; base=h2

\section{第一章}

我\Emph{是}阿黎玛特雅人若瑟，我从比拉多那里求得了主耶稣的身体埋葬，为此我被凶残而敌天主的犹太人关押在监狱中；他们自己也藉着梅瑟遵守律法，共受苦难，激怒了他们的立法者，却不知祂是天主，把祂钉死，并向认识天主的人显示了祂。在那他们判决天主之子被钉死的日子里，在基督受难前七日，两个被定罪的强盗从耶利哥被送到总督比拉多那里；他们的案情如下：——\par

第一个，名叫革斯塔斯，杀害旅客，用刀杀死他们，又把其他人赤身暴露。他把妇女倒挂，头朝下，割下她们的乳房，饮婴儿四肢的血，从不认识天主，不遵法律，自始至终凶暴，行这种事。\par

另一个的案情如下：他名叫德玛斯，是加里肋亚人，开了一家客店。他袭击富人，却善待穷人——一个像多俾亚一样的贼，因为他埋葬穷人的尸体。\BibleRef{多 1:17}他又从事抢劫大批犹太人，在耶路撒冷偷了律法本身，剥光了盖法的女儿（她是圣所的女司祭）的衣服，取走了撒罗满安放在那里的神秘宝物。他的所作所为就是如此。\par

耶稣也是在逾越节前第三天晚上被捉拿。对于盖法和犹太人群众来说，那不是逾越节，而是大哀悼，因为圣所被强盗劫掠。他们召来犹达斯依斯加略，对他说话，因为他是盖法司祭的兄弟的\Emph{儿子}。他并不是在耶稣面前的门徒；但全体犹太人狡猾地支持他，让他跟随耶稣，不是要服从祂所行的奇迹，也不是要承认祂，而是要把祂出卖给他们，希望抓住祂的某个谎言，因这种勇敢而诚实的行径，每天给他半个金协刻耳的礼物。他跟随耶稣做了两年这事，如祂的一个名叫若望的门徒所说。\par

在第三天，耶稣被捉拿前，犹达斯对犹太人说：\Quote{来，我们开个会；因为偷律法的也许不是那强盗，而是耶稣本人，我控告祂。}这些话说了以后，掌管圣所钥匙的尼苛德摩到我们这里来，对众人说：\Quote{不要做这样的事。}因为尼苛德摩是真诚的，超过所有犹太人群众。盖法的女儿，名叫撒辣，喊道说：\Quote{祂自己在众人面前说过反对这圣地的话：\InnerQuote{我能拆毁这座圣殿，并在三天内把它重建起来。}}犹太人对她说：\Quote{你在我们众人面前有价值。}因为他们视她为女先知。确实，会议结束后，耶稣就被捉拿了。\par

% structure: p0009 source=h2 target=section reason=relative-heading-rank; base=h2

\section{第二章}

次日，一周第四天，在第九时辰，他们把祂带入盖法的厅堂。亚纳斯和盖法对祂说：\Quote{告诉我们，你为什么偷了我们的律法，否定了梅瑟和先知的条例？}耶稣一言不答。第二次，群众也在场，他们对祂说：\Quote{撒罗满用四十六年建造的圣所，你为什么想在一瞬间拆毁？}对于这些事，耶稣一言不答。因为会堂的圣所已被那强盗劫掠。\par

第四天的晚上结束时，全体群众想烧死盖法的女儿，因为律法的丢失；因为他们不知道如何守逾越节。她对他们说：\Quote{稍等，我的孩子们，让我们除掉这个耶稣，律法就会找到，圣节就会圆满完成。}亚纳斯和盖法秘密地给了犹达斯依斯加略不少钱，说：\Quote{照你以前对我们说的那样说：\InnerQuote{我知道律法被耶稣偷了}，这样指控就会转向祂，而不会针对这个无罪的少女。}犹达斯接受了这个命令，对他们说：\Quote{不要让全体群众知道我被你们指使做这事来反对耶稣；但释放耶稣，我来说服群众事情就是这样。}他们就狡猾地释放了耶稣。\par

犹达斯在第五天黎明进入圣所，对全体民众说：\Quote{你们愿意给我什么，我就把推翻律法、劫掠先知的人交给你们？}犹太人对他说：\Quote{你若把他交给我们，我们就给你三十块金币。}民众不知道犹达斯说的是耶稣，因为他们中许多人承认祂是天主子。犹达斯就收了那三十块金币。\par

在第四时辰和第五时辰出去，他看见耶稣在街上行走。到了傍晚，犹达斯对犹太人说：\Quote{给我派一些带刀剑棍棒的兵丁相助，我就把他交给你们。}他们就给了他差役去捉拿祂。他们走的时候，犹达斯对他们说：\Quote{我亲吻的那个人，你们就抓住他，因为他偷了律法和先知。}于是他上前亲吻耶稣，说：\Quote{拉比，你好！}那时是第五天的晚上。他们抓住祂，把祂交给盖法和司祭长，犹达斯说：\Quote{这就是偷律法和先知的人。}犹太人给了耶稣不公的审判，说：\Quote{你为什么做这些事？}祂一言不答。\par

尼苛德摩和我若瑟，看见那灾祸的座位，就远离他们，不愿与不敬神者的计谋一同灭亡。\par

% structure: p0015 source=h2 target=section reason=relative-heading-rank; base=h2

\section{第三章}

因此，那一夜他们对耶稣做了许多可怕的事之后，在预备日的黎明，他们把祂交给总督比拉多，让他把祂钉十字架；为此他们都聚集在一起。因此，经过审判，总督比拉多下令把祂与两个强盗一同钉在十字架上。他们与耶稣一同被钉，革斯塔斯在左边，德玛斯在右边。\par

左边的那一个开始喊叫，对耶稣说：\Quote{看，我在世上行了多少恶事；若我知道你是王，我也必砍断你。你为何自称天主子，在急需时却不能救助自己？你怎能将救助赐予另一个祈求帮助的人？你若真是基督，就从十字架上下来，好让我信你。但我现在看见你与我一同丧亡，不像一个人，倒像一头野兽。}他又开始说了许多话攻击耶稣，亵渎他，向他咬牙切齿。原来那强盗活活地被魔鬼的罗网捉住了。\BibleRef{弟后 2:26}\par

但在右边的那强盗，名叫德玛斯，看见耶稣神性的恩宠，就这样喊说：\Quote{我知道你，耶稣基督，你是天主子。我看见你，基督，受到千万天使的朝拜。求你赦免我所犯的罪。在你审判全世界时，不要在我的审判中使星辰或月亮来反对我，因为我在夜间行了我的邪恶计划。不要催促那现在因你而变暗的太阳讲述我心中的恶事，因为我不能给你任何礼物以赦免我的罪。死亡已因我的罪而临于我；但你的赎罪祭是我的。万有的主啊，求你从你可怕的审判中拯救我。不要给仇敌力量吞没我，使他成为我灵魂的继承者，如同那挂在左边的人的灵魂一般；因为我看见魔鬼怎样欢喜地取走他的灵魂，他的身体就消失了。也不要命我往犹太人的分地去；因为我看见梅瑟和先祖们在痛哭，而魔鬼为他们欢喜。所以，主啊，在我灵魂离去之前，求你命我的罪被洗净，并记住我这个罪人，在你的王国里，当你坐在那伟大至高宝座上审判以色列十二支派时。\BibleRef{玛 19:28} 因为你为你自己为你的世界预备了重大的惩罚。}\par

那强盗说完这话，耶稣对他说：\Quote{阿们，阿们，我告诉你，德玛斯，今天你就要同我在乐园里。\BibleRef{路 23:43} 但国度之子，即亚巴郎、依撒格、雅各伯和梅瑟的儿女，将被驱逐到外面的黑暗里；那里要有哀号和切齿。\BibleRef{玛 8:11} 而你，唯独你，要住在乐园里，直到我第二次显现，那时我要审判那些不承认我名的人。}他又对那强盗说：\Quote{你去，告诉革鲁宾和那旋转火焰剑的掌权者，他们自亚当——第一个受造者——在乐园里犯罪、不守我的诫命、被我逐出那时起，就守护着乐园。没有一个先前的人能看见乐园，直到我第二次来临审判生者死者。}他这样写道：\Quote{耶稣基督，天主子，从诸天高处降下，从不可见之父的怀中出来而不与他分离，降世成为血肉，被钉在十字架上，为要拯救我所塑造的亚当——致我的总领天使掌权者，乐园的门卫，我父的官员：我旨意并命令那与我一同被钉的人，应进入乐园，藉我获得罪的赦免；并使他穿上不朽的身体，进入乐园，居住在那无人曾能居住的地方。}\par

看，他说完这话，耶稣便在预备日第九时辰断了气。遍地都黑暗了；因发生大地震，圣所倒塌，圣殿的翼也倒塌了。\par

% structure: p0021 source=h2 target=section reason=relative-heading-rank; base=h2

\section{第四章。}

我若瑟求得了耶稣的身体，把他安放在一个新坟墓里，那里从未葬过人。至于右边那强盗，他的身体找不到了；但左边那强盗的身体，形状如同龙。\par

我求得了耶稣的身体去埋葬后，犹太人被仇恨和愤怒冲昏了头脑，把我关进监牢，那里是拘禁作恶者的地方。这事发生在我身上是在安息日的黄昏，因此我们的民族违犯了法律。看，就是那个民族，在安息日承受了可怕的磨难。\par

如今，在一周的第一天黄昏，在夜间第五时辰，耶稣来到监牢中我面前，同他一起被钉在右边的那强盗也来了，就是被他送入乐园的那位。屋里有大光。那房子被四角悬挂起来，地方开了，我便出来。那时我首先认出了耶稣，又认出那强盗，他带来一封信给耶稣。我们往加里肋亚去时，有一道大光闪耀，这光不是受造物发出的。那强盗身上也发出一股来自乐园的极大馨香。\par

耶稣在一个地方坐下，这样读了：\Quote{我们，革鲁宾和六翼者，被你的神性命令看守乐园花园的，藉着那按你的安排与你同钉的强盗，作如下声明：当我们看见与你同钉的强盗的钉痕，以及你的神性之信的光辉时，火确实熄灭了，不能承受那钉痕的灿烂；我们便蹲伏，极其恐惧。因为我们听说，那创造天地及整个万有的创造主，从高天降下，为亚当——第一个受造者——的缘故，居住在地的低下之处。当我们看见那无玷的十字架如闪电从那强盗身上闪耀，带着七倍于太阳的光辉时，战栗临到了我们。我们感到下界猛烈震动；有大声，阴府的差役与我们一同说：圣、圣、圣，那在起初居于至高者。掌权者也发出呼声：主啊，你在天上和地上显明了，给世界带来喜乐；而比这更大的礼物，你藉着万世无形之计划，使你自己的肖像脱离了死亡。}\par

% structure: p0026 source=h2 target=section reason=relative-heading-rank; base=h2

\section{第五章。}

我目睹了这些事後，与\Person{耶稣}和那强盗同往\Place{加里肋亚}，\Person{耶稣}变了形貌，不再像被钉十字架以前那样，而是完全成为光；天使们不断侍奉祂，\Person{耶稣}与他们交谈。我与祂同住了三天。没有一位祂的门徒与祂同在，只有那强盗一人。\par

在无酵节中期，祂的门徒\Person{若望}来了，我们不再看见那强盗发生何事。\Person{若望}问\Person{耶稣}说：\Quote{这个是谁，祢为何不让我被他看见？}但\Person{耶稣}什么也没有回答他。他俯伏在祂面前说：\Quote{主，我知道祢从起初就爱了我，为何不把那人的事启示给我呢？}\Person{耶稣}对他说：\Quote{你为何寻求那隐藏的事？你仍然不明白吗？你没有觉察到充满这地方的乐园芬芳吗？你不知道他是谁吗？十字架上的强盗成了乐园的继承人。阿们，阿们，我告诉你们：那乐园只属于他一人，直到那大日子的来临。}\Person{若望}说：\Quote{求祢使我堪当看见他。}\par

\Person{若望}还在说话时，那强盗忽然显现；\Person{若望}大惊，扑倒在地。那强盗不再如\Person{若望}来以前那样的最初形态，而是像一位大有权能的君王，身上背着十字架。有一大群人的声音发出说：\Quote{你来到了为你预备好的乐园住处。我们受那派遣你者的命令，要服事你直到那大日子。}这声音过後，那强盗和我\Person{若瑟}都消失了，我被发现在自己家中；此後我不再见\Person{耶稣}。\par

我既目睹这些事，便记录下来，为使众人相信被钉十字架的\Person{耶稣基督}我们的主，不再遵行\Person{梅瑟}的法律，而相信藉着祂所发生的征兆和奇事，并使我们这些相信的人得以承受永生，被寻见于天国。因为荣耀、力量、颂赞、尊威都归于祂，至於无穷之世。阿们。\par

\SourceRecord{Translated by Alexander Walker.；Ante-Nicene Fathers；Vol. 8.；Edited by Alexander Roberts, James Donaldson, and A. Cleveland Coxe.；Buffalo, NY: Christian Literature Publishing Co.,；1886.；<http://www.newadvent.org/fathers/0813.htm>.}
